\vspace{-2mm}
\section{Related Work}
\vspace{-1mm}
\xhdr{Causal studies under network interference.} There have been many causal studies \cite{aronow2017estimating,basse2018analyzing,imai2020causal,kohavi2013online,tchetgen2012causal,ugander2013graph,yuan2021causal,ma2021causal,bhattacharya2020causal} which address the existence of network interference. These works mainly include the following categories: (i) \textit{Random assignment strategy under interference} \cite{aronow2017estimating,basse2018analyzing,imai2020causal,ugander2013graph,fatemi2020minimizing}. These works focus on experimental studies under interference (without SUTVA assumption). %Some of them are based on cluster random assignment [], which treats instances at a group level. 
In some studies \cite{karrer2021network}, strong interference is assumed to exist within each group while there is no interference across different groups; (ii) \textit{Causal effect estimation on observational data with interference} \cite{rakesh2018linked,ma2021causal,arbour2016inferring,tchetgen2021auto}. Different from the experimental studies which can design assignment strategy, another line of works (and also our work) assume interferences exist across individuals in the observational data. They relax the SUTVA assumption and define the potential outcome with a function that takes the instance covariates and treatment assignment of each individual and other interacted individuals as input. Among them, Rakesh et al. \cite{rakesh2018linked} propose a Linked
Causal Variational Autoencoder (LCVA) framework to estimate the causal effect of a treatment on an outcome with the existence of interference between pairs of instances. Different from these works that focus on pairwise spillover effects, Ma et al. \cite{ma2021causal} consider the spillover effect in network structure, and propose a graph neural network (GNN) \cite{kipf2016semi} based framework for causal effect estimation under network interference. %GNNs are used as effective tools to capture the dependencies between nodes and links in the given graph. 
However, these works are still limited in pairs of individuals or ordinary graphs and lack consideration of high-order interference.
{
Another line of studies is bipartite causal inference  \cite{zigler2021bipartite,pouget2019variance}. Traditionally, bipartite causal inference involves two types of units: interventional/outcome units. Interventional units are assigned with treatments, and outcomes are observed from outcome units.  
%In this setting, interference means the outcome of an outcome unit is influenced by multiple interventional units. 
Although this setup is different from ours, considering that there is a node-hyperedge bipartite corresponding to each hypergraph (and ordinary graph), thus %bipartite can also fit in some scenarios of hypergraph. Although 
the two modeling approaches (bipartite and hypergraph) are conceptually similar. Nevertheless, % each has its own advantages in practice. 
we argue hypergraph is a more appropriate framing in many scenarios since: i) hypergraph does not require instantiating edges as additional nodes, or treating these two kinds of nodes differently, thus more computationally efficient; ii) hypergraph has the potential to be more convenient and efficient when generalizing to new hyperedges, while bipartite needs to generate both new nodes for the new hyperedges and their associated new edges.
}

\xhdr{Hypergraph algorithms and neural networks.} To process hypergraph structures for downstream tasks, a line of works simplify the hypergraph structure by taking abstract representations of complicated multi-way interactions \cite{benson2018simplicial,li2013link,xu2013hyperlink,yadati2018link,zhang2018beyond}. Other works directly tackle the original hypergraph structure \cite{arya2018exploiting,benson2018sequences,feng2019hypergraph,huang2015scalable,sharma2014predicting,yadati2018hypergcn}.  
Recently, numerous works have studied on hypergraph neural networks \cite{bai2021hypergraph,yadati2018hypergcn}. Feng el al. \cite{feng2019hypergraph} propose hypergraph neural networks (HGNN) framework to encode high-order data correlation in a hypergraph structure. A hyperedge convolution operation is designed for representation learning. Bai et al. \cite{bai2021hypergraph} introduce two end-to-end trainable operators hypergraph convolution and hypergraph attention to learn node representations in hypergraphs. Yadati et al. \cite{yadati2018hypergcn} develop a self-attention based hypergraph neural network Hyper-SAGNN, which is applicable to homogeneous or heterogeneous hypergraphs with variable hyperedge sizes. Jiang et al. \cite{jiang2019dynamic} propose a dynamic hypergraph neural network (DHGNN) which can dynamically update hypergraph structure on each layer. 